%% kapa: exam=1 round=35 year=2024 period=2 form=A subject=accounting subject_ko=회계학 seq=25 number=65
%% source: Q-Net 감정평가사 기출문제 (cst00309) artlSeq=5213229 file=raw/1008/5213229/2교시 문제지 원본.pdf
%% conversion: kordoc 4.24.51 -> tools/gen1.py (auto); review=proofread-sonnet; flags=tab
\begin{question}{65}
20×1년 1월 1일에 설립된 (주)감평은 확정급여제도를 운영하고 있다. 20×1년도 관련 자료가 다음과 같을 때, 20×1년 말 재무상태표의 기타포괄손익누계액에 미치는 영향은? (단, 확정급여채무 계산 시 적용하는 할인율은 연 10\%이다.)
\begin{framedtext}
\begin{tabular}{@{}p{0.75\linewidth}r@{}}
기초 확정급여채무의 현재가치 & ￦120,000 \tabularnewline
기초 사외적립자산의 공정가치 & 90,000 \tabularnewline
퇴직급여 지급액(사외적립자산에서 기말 지급) & 10,000 \tabularnewline
당기 근무원가 & 60,000 \tabularnewline
사외적립자산에 기여금 출연(기말 납부) & 20,000 \tabularnewline
기말 확정급여채무의 현재가치 & 190,000 \tabularnewline
기말 사외적립자산의 공정가치 & 110,000 \tabularnewline
\end{tabular}
\end{framedtext}
\begin{choices}
  \item ￦2,000 감소
  \item ￦2,000 증가
  \item 영향없음
  \item ￦7,000 감소
  \item ￦7,000 증가
\end{choices}
\end{question}
\answer{4}
